% 3-7-ml-dl.tex
%  Short intro to the ML/DL subsection -- heading plus framing only;
%  content lives in 3-7-1 (LightGBM) and 3-7-2 (LSTM).
%  Source map: tree ensembles and deep architectures as the two dominant
%   non-linear families in recent EPF work -> jiang_review_2018 ,
%   oconnor_review_2025 ; the deliberately narrow model list and the
%   equal-budget tuning argument are this project's own.
%  Optuna is a proper noun -> no footnote, no expansion (HANDOFF A3).

\section{مدل‌های یادگیری ماشین و یادگیری عمیق}
\label{sec:3-7}

دو مدل غیرخطی در این پژوهش به کار می‌روند: \lr{LightGBM} به‌عنوان نماینده‌ی خانواده‌ی
درختی و \lr{LSTM} به‌عنوان نماینده‌ی خانواده‌ی یادگیری عمیق. این دو، پرکاربردترین
خانواده‌های غیرخطی در ادبیات اخیر پیش‌بینی قیمت برق‌اند
\cite{jiang_review_2018, oconnor_review_2025}.

فهرست مدل‌ها عامدانه کوتاه است. جنگل تصادفی، \lr{XGBoost}، ماشین بردار پشتیبان و
\lr{GRU} همگی بررسی و از دامنه‌ی پژوهش کنار گذاشته شدند. منطق این محدودسازی آن است که
افزودن گونه‌های بیشتر از یک خانواده، به‌جای پاسخ به پرسش «کدام خانواده‌ی مدل در این
بازار بهتر عمل می‌کند»، بودجه‌ی تنظیم و توان تحلیل را میان گونه‌هایی پخش می‌کند که
رفتار مشابهی دارند. یک نماینده‌ی خوب‌تنظیم‌شده از هر خانواده، اطلاعات بیشتری می‌دهد تا
پنج نماینده‌ی کم‌تنظیم.

سه قید بر هر دو مدل این بخش یکسان اعمال می‌شود، و همین یکسانی است که مقایسه‌ی آن‌ها را
معنادار می‌کند.

نخست، \textbf{بودجه‌ی برابر تنظیم}: هر دو مدل با ۵۰ آزمایش \lr{Optuna} تنظیم شده‌اند،
بر پنجره‌ی اعتبارسنجیِ یکسان که اکیداً پیش از دوره‌ی آزمون قرار دارد. این برابری، پاسخ
پیشاپیشِ پرسشی است که به‌درستی پرسیده خواهد شد: آیا مدل عمیق صرفاً کم‌تنظیم نشده
است؟ بودجه یکسان بوده و در متن ثبت شده است.

دوم، \textbf{تابع زیان هم‌تراز}: هر دو مدل بر زیان قدرمطلق آموزش می‌بینند، هم‌راستا
با معیار سرشناسِ گزارش در فصل چهارم. بدون این هم‌ترازی، بخشی از اختلاف میان دو مدل
به تفاوت در تابع زیان بازمی‌گشت و مقایسه به مصنوعی از انتخاب زیان بدل می‌شد.

سوم، \textbf{منبع یکتای ویژگی}: هر دو از همان ماتریس ۲۴۷ ستونیِ \cref{sec:3-4} تغذیه
می‌کنند و هیچ‌یک مسیر ویژگی‌سازیِ اختصاصی ندارد.

\cref{sec:3-7-1} مدل درختی و \cref{sec:3-7-2} مدل عمیق را شرح می‌دهد. در هر دو، آهنگ واسنجی، رویه‌ی
تنظیم و سازوکارهای تضمین بازتولیدپذیری به‌تفصیل می‌آید، چون همین جزئیات‌اند که
تعیین می‌کنند اعداد فصل چهارم تکرارپذیر باشند یا نه.
